\newpage
\section{Empfangen von FSK-Modulierten Signalen}\label{sec:task1}

% \begin{itemize}[noitemsep]
%     \item Kurze Beschreibung der Aufgabenstellung
%     \item Screenshot des GNU Radio Companion Aufbaus
%     \item Berechnung der Frequenzabweichung
%     \item Analyse des empfangenen Signals im Zeitbereich
%     \begin{itemize}[noitemsep]
%         \item Grafische Darstellung
%         \item Interpretation des Ergebnisses (z.B. stimmt die berechnete / eingestellte Frequenzabweichung mit der Messung überein)
%         \item Vergleich mit der Sequenz aus dem Datenblatt (Abb. 2 auf Seite 42)
%     \end{itemize}
%     \item Analyse des empfangenen Signals im Frequenzbereich
%     \begin{itemize}[noitemsep]
%         \item Grafische Darstellung
%         \item Interpretation des Ergebnisses
%     \end{itemize}
% \end{itemize}

Ein im Labor platziertes eLITe-Board sendet dauerhaft ein 2-FSK moduliertes Signal aus. Dieses Signal soll nun mit dem USRP empfangen werden und in GNU Radio Companion (GRC) demoduliert und ausgewertet werden. Die Einstellungen des Senders sind dabei:

\begin{itemize}[noitemsep]
    \item $f_\text{osc} = \SI{26}{\mega\hertz}$
    \item Baudrate: $\SI{1}{\kilo\bit\per\second}$
    \item Sync Word: \texttt{0x0F0F}
    \item Frequenzabweichung: Exponent \texttt{0b100}, Mantisse \texttt{0b111} 
    \item Data Word: \texttt{0xA5}
\end{itemize}

\subsection{Berechnung der Frequenzabweichung}

Gegeben sind nur die Registerwerte der Konfiguration des CC1101. Die Frequenzabweichung kann mit folgender Formel berechnet werden \cite[S.~42]{TI_CC1101_Datasheet}:
\begin{equation}
    f_\text{dev} = \frac{f_\text{xosc}}{2^{17}} (8+\text{DEVIATION\_M})\cdot 2^{\text{DEVIATION\_E}}
\end{equation}
Mit den Sender-Einstellungen aus \autoref{sec:task1} ist die Frequenzabweichung:
\begin{equation}
    f_\text{dev} = \frac{\SI{26}{\mega\hertz}}{2^{17}} (8+7)\cdot 2^{4} = \SI{47607}{\hertz}\label{eq:f_dev}
\end{equation}

\subsection{Aufbau in GRC}

Zur demodulation des empfangenen Signals wurde der in \autoref{fig:task1_grc} Aufbau in GNU Radio Companion verwendet. Die zuvor berechnete Frequenzabweichung wird in einer Variable gespeichert, um sie im \textit{Quadrature Demod} Block zu verwenden.

\begin{figure}[h]
    \centering
    \includegraphics[width=\textwidth]{\usrpRecvPath{1_grc.pdf}}
    \caption{GNU Radio Companion Aufbau für Task 1}\label{fig:task1_grc}
\end{figure}

Um das Signal zu demodulieren wurde der \textit{Quadrature Demod} Block
verwendet. Mit Quadrature Demod wird eine Änderung in der Phase eines
Komplexen Samples in einen Proportionalen Floatingpoint-Wert umgewandelt. Dabei wird der
komplexe Abtastwert mit dem konjugiert komplexen vorherigen Abtastwert
multipliziert, und die Phase Ausgelesen. So kann die Informationstragende Phase
bzw. Phasenänderung des FSK-Signals wieder extrahiert werden.

Mit einem \textit{Rational Resampler} wird die Abtastrate des Signals um den Faktor $100$ heruntergesetzt, um das Signal in einem größeren Zeitfenster betrachten zu können. Hier ist zu beachten, dass die Abtastrate des \textit{QT GUI Time Sinks} um den selben Faktor reduziert werden muss, damit die Zeitachse korrekt dargestellt wird.

\subsection{Analyse im Zeitbereich}

Das demodulierte Signal im Zeitbereich ist in \autoref{fig:task1_time} dargestellt. Hier ist das Bitmuster der gesendeten Daten zu erkennen. 

\begin{figure}[h]
    \begin{minipage}[t]{.5\textwidth}
    \centering
    \includegraphics[height=4cm]{\usrpRecvPath{1_Bit_correct_time_scale_annotated.png}}
    \caption{Demoduliertes Zeitsignal}\label{fig:task1_time}
    \end{minipage}\hfill
    \begin{minipage}[t]{.5\textwidth}
    \centering
    \includegraphics[height=4cm]{\usrpRecvPath{1_Bit_Baudrate.png}}
    \caption{Baudrate des empfangenen Signals}\label{fig:task1_baudrate}
    \end{minipage}
\end{figure}

Aus dem demodulierten Signal können die einzelnen Symbole abgelesen werden. Diese sind in der \autoref{tab:task1_received_data} zusammengefasst.

\begin{table}[h]
    \centering
    \caption{Aus dem Messsignal extrahierte Daten}\label{tab:task1_received_data}
    \begin{tabularx}{\textwidth}{YYY}
        Präambel & Sync Word & Data Word \\ \midrule
        \texttt{0xAAAA AAAA} & \texttt{0x0F0F 0F0F} & \texttt{0xA5} \\
    \end{tabularx}
\end{table}

Ein Vergleich mit dem im Datenblatt angegebenen Frame-Format in \autoref{fig:task1_cc1101_frame} zeigt, dass das empfangene Signal dem erwarteten Muster entspricht. Zu beachten ist, dass das Sync Word  am Sender von \texttt{0xD3} auf \texttt{0x0F} geändert wurde.

\begin{figure}[h]
    \centering
    \includegraphics[width=.8\textwidth]{\usrpRecvPath{cc1101_fig_21.png}}
    \caption{Daten-Format aus dem Datenblatt \cite[S.~42, Abb.~21]{TI_CC1101_Datasheet}}\label{fig:task1_cc1101_frame}
\end{figure}





Anhand der Zeitachse in \autoref{fig:task1_baudrate} ist zu erkennen, dass die Periode eines Bits ca. $\SI{1}{\milli\second}$ beträgt, was mit $\frac{1}{\SI{1}{\kilo\bit\per\second}}$ übereinstimmt und die Baudrate laut Senderkonfiguration bestätigt.



\newpage
\subsection{Analyse im Frequenzbereich}

Im Spektrum des empfangenen Signals sind die beiden Frequenzanteile der FSK-Modulation gut zu erkennen. Der gemessene Frequenzabstand der Peaks zur Mittelfrequenz in \autoref{fig:task1_spectrum1} beträgt ca. $\pm\SI{50}{\kilo\hertz}$, was gut mit der berechneten Abweichung von $\SI{47.6}{\kilo\hertz}$ übereinstimmt.

\begin{figure}[h]
    \centering
    \includegraphics[width=0.7\textwidth]{\usrpRecvPath{1_Spectrum.png}}
    \caption{Empfangenes Signal im Frequenzbereich}\label{fig:task1_spectrum1}
\end{figure}

Im Zeitlichen Verlauf des Spektrums in \autoref{fig:task1_spectrum2} ist zu erkennen, dass die beiden Frequenzanteile des FSK-Signals regelmäßig wechseln, was dem gesendeten Datensignal entspricht. 

\begin{figure}[h]
    \centering
    \includegraphics[width=0.7\textwidth]{\usrpRecvPath{1_Waterfall.png}}
    \caption{Empfangenes Signal im Frequenzbereich}\label{fig:task1_spectrum2}
\end{figure}

Dabei ist die untere Frequenzkomponente dem Symbol \texttt{'0'} und die obere dem Symbol \texttt{'1'} zugeordnet.

\begin{table}[h]
    \centering
    \caption{FSK Symbolencoding \cite[S.~42, Tab.~29]{TI_CC1101_Datasheet}}
    \label{tab:fsk_modulation}
    \begin{tabularx}{\textwidth}{YYY}
        Format & Symbol & Coding \\ \midrule
        \multirow{2}{*}{2-FSK} & \texttt{'0'} & $-f_\text{dev}$ \\
        & \texttt{'1'} & $+f_\text{dev}$ \\
    \end{tabularx}
\end{table}
